Towing kites, wing sails and Flettner rotors all sail on cargo ships today, but their published savings come from different ships, routes and methods, so nobody can say which rig is better where. We compare the three with one model on the same ships and routes. Each ship is weather-routed through three years of real Atlantic weather and chooses its own speed, and each rig is sized for the lowest total cost per tonne-kilometre, which includes the rig's price, its wear, fuel and the value of time. We do this for thirteen ship classes, from a coaster to the largest container ships, on six routes and three fuels: fuel oil, e-methanol and hydrogen. The kite is a young technology, so we set goals for what development can still change, such as canopy life and price and how often a rig fails, and give each rig's results at its goals. The result reads: if each rig reaches its goals, it saves this much.

On fuel oil no rig saves much, a few percent of the cost, because the fuel is cheap. On e-methanol and hydrogen the savings are larger: kites and rotors each save about \tbd{a tenth} of the cost, level within the accuracy of this draft, and wing sails about half as much. Which rig leads depends on the ship and the route: kites on container ships and on the routes to the Caribbean, the Gulf of Mexico and Bilbao, rotors on the car carrier and on the routes to Dakar and Brazil, with bulk carriers and tankers about level. No rig closes the gap between fuel oil and e-methanol: the best closes about \tbd{a third} of it. Part of the saving is time, since a ship with a rig sails faster at the same cost. The kite gets this far only at its goals: with a canopy that lasts much shorter than the goal, or a kite lost now and then, it falls behind the rotors, which already have a service record. A rigid kite flies at ship scale but does not pay. A companion website lets readers change the assumptions.
